\begin{numlemma}[Bathtub compactness at the radial profile]\label{lem:bathtub-compactness-radial}
Let
\[
w_\alpha(z):=(1-|z|^2)^{\alpha+2},
\qquad
\Theta(s):=\sup_{\tau(E)=s}\int_E w_\alpha\,d\tau
=\int_{B_s}w_\alpha\,d\tau .
\]
If $u_j\to w_\alpha$ in $L^1(d\tau)$ and $\tau(E_j)=s$ with
\[
\int_{E_j}u_j\,d\tau\longrightarrow \Theta(s),
\]
then
\[
\tau(E_j\triangle B_s)\longrightarrow0.
\]
In particular,
\[
\int_{B_s}u_j\,d\tau\longrightarrow\Theta(s).
\]
\end{numlemma}

\begin{proof}
Let $B_s=\{|z|<r_s\}$. Fix $\varepsilon>0$ and choose an annulus
\[
A_\varepsilon:=\{r_s-\varepsilon<|z|<r_s+\varepsilon\}
\]
with $\tau(A_\varepsilon)<\varepsilon$. On
$B_s\setminus A_\varepsilon$ the radial function $w_\alpha$ is strictly larger
than on $\D\setminus(B_s\cup A_\varepsilon)$; denote the gap by
\[
\gamma_\varepsilon
:=
\inf_{B_s\setminus A_\varepsilon}w_\alpha
-
\sup_{\D\setminus(B_s\cup A_\varepsilon)}w_\alpha
>0.
\]
Since $\tau(E_j)=\tau(B_s)$, the mass of $B_s\setminus E_j$ must be replaced
by the same amount of mass in $E_j\setminus B_s$, up to the annular error.
Therefore the strict rearrangement inequality gives, after possibly reducing
the constant,
\[
\Theta(s)-\int_{E_j}w_\alpha\,d\tau
\ge
\gamma_\varepsilon
\tau\bigl((E_j\triangle B_s)\setminus A_\varepsilon\bigr)
-\gamma_\varepsilon\tau(A_\varepsilon).
\]
The same deficit with $u_j$ differs from this one by at most
$2\|u_j-w_\alpha\|_{L^1(d\tau)}$, and therefore
\[
\limsup_j
\tau\bigl((E_j\triangle B_s)\setminus A_\varepsilon\bigr)=0.
\]
Letting $\varepsilon\downarrow0$ proves the symmetric-difference convergence.
The final assertion follows from
\[
\left|\int_{B_s}u_j\,d\tau-\Theta(s)\right|
\le \|u_j-w_\alpha\|_{L^1(d\tau)}.
\]
\end{proof}

\begin{numlemma}[Harmonic bathtub Schur complement]\label{lem:harmonic-bathtub-schur}
Fix $\alpha>-1$ and $s\in(0,\infty)$, and let
$B_s=\{|z|<r_s\}$, $R:=r_s^2$, $\theta:=\theta_\alpha(s)$. For a real-valued
harmonic function
\[
h(re^{i\vartheta})
=h(0)+\sum_{k\ge1} r^k
\Re(b_k e^{ik\vartheta})
\]
on a neighbourhood of $\overline\D$, set
\[
g_h:=\frac{e^h}{\|e^h\|_{2,\alpha}},
\qquad
\mathcal D_s(h):=
\theta-\sup_{\tau(E)=s}\int_E g_h^2\,dA_\alpha .
\]
There are $\eta_{\alpha,s}>0$ and $c_{\alpha,s}>0$ such that, if
$\|h-h(0)\|_{C^2(\overline\D)}\le\eta_{\alpha,s}$, then
\[
\mathcal D_s(h)
\ge
c_{\alpha,s}\sum_{k\ge2}\|r^k\Re(b_k e^{ik\vartheta})\|_{2,\alpha}^2 .
\]
Equivalently,
\[
\inf_{|a|<1}
\left\|
\frac{e^h}{\|e^h\|_{2,\alpha}}-g_a
\right\|_{2,\alpha}^2
\le
C_{\alpha,s}\,\mathcal D_s(h).
\]
\end{numlemma}

\begin{proof}
Constants in $h$ disappear from $g_h$, so assume $h(0)=0$. The bathtub
maximizer for $g_h^2dA_\alpha$ under the constraint $\tau(E)=s$ is a superlevel
set of
\[
e^{2h(z)}(1-|z|^2)^{\alpha+2}.
\]
Since
\[
\partial_r(1-r^2)^{\alpha+2}\big|_{r=r_s}
=-2(\alpha+2)r_s(1-r_s^2)^{\alpha+1}<0,
\]
the implicit function theorem gives, for $\|h\|_{C^2}$ small, a boundary graph
\[
E_h=\{re^{i\vartheta}:0<r<r_s+\rho(\vartheta)\},
\qquad
\int_0^{2\pi}\rho(\vartheta)\,d\vartheta=O(\|\rho\|_{L^2}^2),
\]
for the maximizing set. The last relation is the first-order form of the
constraint $\tau(E_h)=\tau(B_s)$.

It remains to compute the quadratic part of
$\theta-\int_{E_h}g_h^2\,dA_\alpha$. Write
\[
dA_\alpha=A(r)\,dr\,d\vartheta,\qquad
d\tau=T(r)\,dr\,d\vartheta,
\]
where
\[
A(r)=\frac{\alpha+1}{\pi}(1-r^2)^\alpha r,
\qquad
T(r)=\frac{r}{(1-r^2)^2}.
\]
The Lagrange multiplier for the unperturbed bathtub problem is
\[
\lambda_s:=\frac{A(r_s)}{T(r_s)}
=\frac{\alpha+1}{\pi}(1-r_s^2)^{\alpha+2}.
\]
For variations preserving $\tau$ to first order, the second-order expansion of
\[
\int_E g_h^2\,dA_\alpha-\lambda_s\tau(E)
\]
at $(h,E)=(0,B_s)$ is diagonal in angular Fourier modes. The block associated
to the $k$th harmonic coefficient and the $k$th boundary displacement is
\[
\Gamma_k |b_k|^2
-2A(r_s)r_s^k
\int_0^{2\pi}\Re(b_ke^{ik\vartheta})\,\rho_k(\vartheta)\,d\vartheta
+\kappa_s\|\rho_k\|_{L^2(\partial\D)}^2,
\]
where
\[
\kappa_s
:=
-\frac12 T(r_s)\left(\frac{A}{T}\right)'(r_s)
=
\frac{(\alpha+1)(\alpha+2)}{\pi}
r_s^2(1-r_s^2)^{\alpha-1}>0,
\]
and
\[
\Gamma_k
:=
\theta M_k-I_k(R),
\quad
M_k:=(\alpha+1)\int_0^1 t^k(1-t)^\alpha\,dt,
\quad
I_k(R):=(\alpha+1)\int_0^R t^k(1-t)^\alpha\,dt .
\]
Optimizing in the boundary coefficient gives the Schur complement
\[
\mathfrak s_k
=
\Gamma_k
-
\frac{\alpha+1}{\alpha+2}(1-R)^{\alpha+1}R^k .
\]
The cancellation $\mathfrak s_1=0$ is exactly the infinitesimal Mobius
translation. For $k\ge2$ the coefficient is strictly positive. Indeed, with
$d\nu(t)=(\alpha+1)(1-t)^\alpha\,dt$,
\[
\Gamma_k
=
\int_0^R\int_R^1 (y^k-x^k)\,d\nu(y)d\nu(x),
\]
and the subtracted term is $R^{k-1}$ times the same expression with $k=1$.
Thus
\[
\mathfrak s_k
=
\int_0^R\int_R^1
\bigl(y^k-x^k-R^{k-1}(y-x)\bigr)\,d\nu(y)d\nu(x)>0
\qquad (k\ge2),
\]
because $x<R<y$ and
\[
y^k-x^k=(y-x)\sum_{j=0}^{k-1}y^{k-1-j}x^j
>R^{k-1}(y-x).
\]
Moreover $\mathfrak s_k/M_k$ has a positive lower bound for $k\ge2$: the
displayed strict positivity handles finitely many $k$, while
$\mathfrak s_k/M_k\to\theta$ as $k\to\infty$.

Consequently the quadratic part controls precisely the harmonic modes
$k\ge2$. The cubic remainder is bounded by
$C_{\alpha,s}\|h\|_{C^2}\bigl(\sum_{k\ge2}\|h_k\|_{2,\alpha}^2+
\|\rho\|_{L^2}^2\bigr)$, and is absorbed after decreasing
$\eta_{\alpha,s}$. This proves the first estimate. The second estimate follows
from the tubular expansion of the extremal manifold at $g_0=1$:
the $k=1$ harmonic modes are its tangent space, and the modes $k\ge2$ are
transversal.
\end{proof}
